\documentclass[10pt]{beamer}

\usepackage{appendixnumberbeamer}
\usepackage{booktabs}
\usepackage[scale=2]{ccicons}
\usepackage{pgfplots}
\usepackage{xspace}
\usepackage{bookmark}
\usepackage{amssymb}
\usepackage{mathtools}
\usepackage[normalem]{ulem}
\usepackage[T1]{fontenc}
\usepackage[sfdefault,book]{FiraSans}
\usepackage{FiraMono}
\usepackage{fontawesome}
\usepackage{hyperref}
\usepackage{listings}

\renewcommand{\thefootnote}{\fnsymbol{footnote}}
\renewcommand{\MintedPygmentize}{/home/thiago/.local/bin/pygmentize}

\definecolor{light-gray}{gray}{0.95}
\renewcommand\lstlistingname{Código}
\DeclareCaptionFormat{listing} {
  \parbox{\textwidth}{\hspace{-0.2cm}#1#2#3}
}
\DeclareCaptionFont{black}{\color{black}}
\captionsetup[lstlisting]{
  format=listing,
  labelfont=black,
  textfont=black,
  singlelinecheck=true,
  margin=0pt,
  font={tt,footnotesize,bf}
}

\lstset{
  basicstyle=\footnotesize\ttfamily,
  escapeinside={\%*}{*)},
  mathescape=true,
  showspaces=false,
  showtabs=false,
  showstringspaces=false,%
  rulesepcolor=\color{black},
  frame=shadowbox,
  upquote=true,
  literate=
  {á}{{\'a}}1 {é}{{\'e}}1 {í}{{\'i}}1 {ó}{{\'o}}1 {ú}{{\'u}}1
  {Á}{{\'A}}1 {É}{{\'E}}1 {Í}{{\'I}}1 {Ó}{{\'O}}1 {Ú}{{\'U}}1
  {à}{{\`a}}1 {è}{{\`e}}1 {ì}{{\`i}}1 {ò}{{\`o}}1 {ù}{{\`u}}1
  {À}{{\`A}}1 {È}{{\'E}}1 {Ì}{{\`I}}1 {Ò}{{\`O}}1 {Ù}{{\`U}}1
  {ä}{{\"a}}1 {ë}{{\"e}}1 {ï}{{\"i}}1 {ö}{{\"o}}1 {ü}{{\"u}}1
  {Ä}{{\"A}}1 {Ë}{{\"E}}1 {Ï}{{\"I}}1 {Ö}{{\"O}}1 {Ü}{{\"U}}1
  {â}{{\^a}}1 {ê}{{\^e}}1 {î}{{\^i}}1 {ô}{{\^o}}1 {û}{{\^u}}1
  {Â}{{\^A}}1 {Ê}{{\^E}}1 {Î}{{\^I}}1 {Ô}{{\^O}}1 {Û}{{\^U}}1
  {Ã}{{\~A}}1 {ã}{{\~a}}1 {Õ}{{\~O}}1 {õ}{{\~o}}1
  {œ}{{\oe}}1 {Œ}{{\OE}}1 {æ}{{\ae}}1 {Æ}{{\AE}}1 {ß}{{\ss}}1
  {ű}{{\H{u}}}1 {Ű}{{\H{U}}}1 {ő}{{\H{o}}}1 {Ő}{{\H{O}}}1
  {ç}{{\c c}}1 {Ç}{{\c C}}1 {ø}{{\o}}1 {å}{{\r a}}1 {Å}{{\r A}}1
  {€}{{\euro}}1 {£}{{\pounds}}1 {«}{{\guillemotleft}}1
  {»}{{\guillemotright}}1 {ñ}{{\~n}}1 {Ñ}{{\~N}}1 {¿}{{?`}}1
}

\setminted{
  fontsize=\footnotesize,
  style=bw,
  frame=single,
  labelposition=topline,
}

\pretitle{\begin{center}\normalsize\bfseries\MakeUppercase{Programação 2 -- }\MakeUppercase}
\author{\normalsize Prof. Thiago Cavalcante}
\date{\vspace{-4ex}}

\geometry{
  top=0.0cm,
  left=1.0cm,
  right=1.0cm,
  bottom=1.0cm
}

\setlength\columnsep{30pt}

\usetikzlibrary{arrows}

\makeatletter
\AddEnumerateCounter{\PaddingUp}{\two@digits}{A00}
\AddEnumerateCounter{\PaddingDown}{\two@digits}{A00}
\newcommand\PaddingUp[1]{\expandafter\two@digits\csname c@#1\endcsname}
\newcommand\PaddingDown[1]{\PaddingUp{#1}\addtocounter{#1}{-2}}
\makeatother

\makeatletter
\def\enumalphalphcnt#1{\expandafter\@enumalphalphcnt\csname c@#1\endcsname}
\def\@enumalphalphcnt#1{\AlphAlph{#1}}
\makeatother
\AddEnumerateCounter{\enumalphalphcnt}{\@enumalphalphcnt}{aa}

\pagenumbering{gobble}

\title{Programação 2}
\author{Thiago Cavalcante  -- thiago.cavalcante@penedo.ufal.br}
\institute{Universidade Federal de Alagoas -- UFAL \\ Campus Arapiraca \\ Unidade de Ensino de Penedo}
\titlegraphic{\hfill\includegraphics[height=1.5cm]{images/brasao-ufal.eps}}


\subtitle{Aula 1}
\date{23 de outubro de 2019}

\begin{document}

\maketitle

\begin{frame}{Sumário}
  \setbeamertemplate{section in toc}[sections numbered]
  \setbeamertemplate{subsection in toc}[ball unnumbered]
  \tableofcontents
\end{frame}

\section{Introdução}

\subsection{Informações Gerais da Disciplina}

\begin{frame}[fragile]{Informações Gerais da Disciplina}
    \begin{center}
        \begin{tabular}{@{}ll@{}}
            \toprule
            \textbf{Nome} & Programação 2 \\
            \textbf{Código} & SISB087 \\
            \textbf{Semestre} & 3º \\
            \textbf{Carga Horária} & 72h \\
            \textbf{PPC} & 02/2019 \\
            \textbf{Turma} & 2019.2 \\
            \textbf{Horário} & Quarta, 19:00 -- 22:30 \\
            \textbf{Sala} & 5 \\
            \textbf{Lista de discussão} & Google Groups: \href{https://groups.google.com/forum/#!forum/sisb087_20192}{\alert{sisb087\_20192}} \\
            \textbf{Repositório} & GitHub: \href{https://github.com/theagoliveira/sisb087_20192}{\alert{theagoliveira/sisb087\_20192}} \\
            \bottomrule
        \end{tabular}
    \end{center}
\end{frame}

\subsection{Pesquisa}

\begin{frame}{Pesquisa}
    \setbeamertemplate{itemize/enumerate body begin}{\large}

    \begin{enumerate}
    \item Pode trazer um computador para aulas práticas?
    \item Já usou \textbf{Git}?
    \item Já usou \textbf{GitHub}?
    \item Consegue compreender textos em \textbf{inglês}?
    \end{enumerate}

    \bigskip
    \bigskip

    \centering
    \texttt{<NOME COMPLETO> <E-MAIL> <RESPOSTAS>}
\end{frame}

\subsection{Recapitulando: Programação 1}

\begin{frame}{Recapitulando: Programação 1}
    \begin{columns}
    \begin{column}{0.5\textwidth}
        \minipage[c][0.7\textheight][s]{\columnwidth}
            \begin{itemize}
                \item Linguagem: \textbf{C}
                \item Conceitos e técnicas de programação básica, procedimentos, algoritmos e programas
                \item Identificadores, constantes, variáveis e atribuição
                \item Tipos primitivos de dados
                \item Comandos de entrada e saída
                \item Operadores, funções e expressões
                \item Instruções condicionais e de repetição
            \end{itemize}
            \vfill
        \endminipage
    \end{column}
    \begin{column}{0.5\textwidth}
        \minipage[c][0.7\textheight][s]{\columnwidth}
            \begin{itemize}
                \item Tipos definidos pelo programador e tipos abstratos de dados
                \item Noções de ponteiros
                \item Estruturas compostas de dados: vetores, matrizes e registros
                \item Manipulação de uma cadeia de caracteres
                \item Noções de arquivos
                \item Programação de algoritmos usando uma LP estruturada
                \item Boas práticas de programação
            \end{itemize}
            \vfill
        \endminipage
    \end{column}
    \end{columns}
\end{frame}

\subsection{Ementa}

\begin{frame}{Ementa}
    \begin{itemize}
        \item Linguagem: \textbf{C}
        \item Importância das estruturas de dados na solução de problemas
        \only<1>{\item Vetores e matrizes}
        \only<2>{\item \sout{Vetores e matrizes}}
        \item Estruturas de dados lineares e não lineares
        \item Pilhas, filas, listas, árvores, florestas, introdução à grafos
        \item Implementação de estruturas de dados com alocação estática e dinâmica de memória
        \item Implementação de estruturas de dados com e sem ponteiros
        \item Algoritmos de ordenação
        \item Algoritmos de busca
        \item Programação avançada e resolução de problemas complexos
        \item Introdução à análise de algoritmos
    \end{itemize}
\end{frame}

\subsection{Bibliografia}

\begin{frame}{Bibliografia}
    \huge
    Livros sugeridos no \alert{repositório}

    \Large
    \begin{itemize}
        \item Programação em C
        \only<1>{\item Estruturas de dados com exemplos em C}
        \only<2>{\item \alert{Estruturas de dados com exemplos em C}}
        \item Estruturas de dados com exemplos em outras linguagens
    \end{itemize}
\end{frame}

\subsection{Cronograma}

\begin{frame}{Cronograma}
    \begin{center}
        \begin{tabular}{@{}r@{/}r@{/}rll@{}}
            \toprule
            \multicolumn{4}{@{}l}{\textbf{Datas importantes}} \\
            \midrule
            13 & 11 & 2019 & \textit{Possível} dia da \alert{AB1} \\
            20 & 11 & 2019 & Feriado \\
            23 & 11 & 2019 & Prazo final para digitação da \alert{AB1} \\
            12 & 02 & 2020 & \textit{Possível} dia da \alert{AB2} \\
            17 & 02 & 2020 & Prazo final para digitação da \alert{AB2} \\
            17--22 & 02 & 2020 & Período de \alert{reavaliação} \\
            27--29 & 02 & 2020 & Período de \alert{provas finais} \\
            \bottomrule
        \end{tabular}
    \end{center}
\end{frame}

\section{Estruturas de Dados}

\setbeamertemplate{frame footer}{SCHILDT, H. C Completo e Total}
\begin{frame}{Introdução a Estruturas de Dados}
    \Large

    \onslide<1->{
    \begin{center}
        \setlength{\fboxsep}{2\fboxsep}\boxed{\textbf{$\text{Programa} = \text{Algoritmos} + \text{Estruturas de Dados}$}}
    \end{center}
    }

    \bigskip

    \onslide<2->{
    Escolha e implementação da estrutura de dados \textbf{são tão importantes quanto} as rotinas
    }

    \bigskip

    \onslide<3->{
    O tipo de estrutura é \textbf{determinado pela natureza} do problema
    }
\end{frame}
\setbeamertemplate{frame footer}{}

\setbeamertemplate{frame footer}{SKIENA, S. S. The Algorithm Design Manual}
\begin{frame}{Introdução a Estruturas de Dados}
    \Large

    \onslide<1->{
    \textit{Tipos abstratos de dados} podem ser implementados \textbf{corretamente} com diferentes estruturas de dados
    }

    \bigskip

    \onslide<2->{
    A forma como as operações são realizadas por uma estrutura podem \textbf{melhorar} (ou \textbf{piorar}) drasticamente a performance de um programa
    }

    \bigskip

    \onslide<3->{
    Projeto: melhor prevenir do que remediar
    }
\end{frame}
\setbeamertemplate{frame footer}{}

\setbeamertemplate{frame footer}{LOUDON, K. Mastering Algorithms with C}
\begin{frame}{Introdução a Estruturas de Dados}
    \Large

    Três motivos para usar estruturas de dados:

    \large

    \begin{itemize}
        \setlength\itemsep{1em}
        \onslide<1->{\item \textbf{Eficiência}: a forma de organização dos dados pode aumentar a \alert{velocidade de execução} de algoritmos (de busca e/ou ordenação, por exemplo)}
        \onslide<2->{\item \textbf{Abstração}: melhor \alert{entendimento} dos dados na solução de problemas}
        \onslide<3->{\item \textbf{Reusabilidade}: estruturas de dados são \alert{modulares} e \alert{independentes de contexto}}
    \end{itemize}
\end{frame}
\setbeamertemplate{frame footer}{}

\begin{frame}{Introdução a Estruturas de Dados}
    \Large

    Características:

    \begin{itemize}
        \onslide<1->{\item Estrutura Linear x Não-linear}
        \onslide<2->{\item Estrutura Contígua x Ligada}
        \onslide<3->{\item Dados Homogêneos x Heterogêneos}
        \onslide<4->{\item Alocação de Memória Estática x Dinâmica}
    \end{itemize}
\end{frame}

\setbeamertemplate{frame footer}{ZIVIANI, N. Projetos de Algoritmos Com Implementações em Pascal e C}
\begin{frame}{Tipos Abstratos de Dados -- TAD}
    \Large

    "Um TAD pode ser visto como um modelo matemático, acompanhado das operações definidas sobre o modelo."
\end{frame}
\setbeamertemplate{frame footer}{}

\setbeamertemplate{frame footer}{SEDGEWICK, R. Algorithms in C, Parts 1-4: Fundamentals, Data Structures, Sorting, Searching}
\begin{frame}{Tipos Abstratos de Dados -- TAD}
    \Large

    "Um TAD é um tipo de dados (conjunto de valores e operações sobre esses valores) que é acessado apenas através de uma interface."
\end{frame}
\setbeamertemplate{frame footer}{}

\setbeamertemplate{frame footer}{CELES, W.; CERQUEIRA, R.; RANGEL, J. L. Introdução a Estruturas de Dados}
\begin{frame}{Tipos Abstratos de Dados -- TAD}
    \Large

    "Nos casos em que um módulo define um novo tipo de dado e o conjunto de operações para manipular dados desse tipo, dizemos que o módulo representa um TAD. Nesse contexto, abstrato significa 'esquecida a forma de implementação', ou seja, um TAD é descrito pela finalidade do tipo e de suas operações, e não pela forma como está implementado."
\end{frame}
\setbeamertemplate{frame footer}{}

\begin{frame}{Tipos Abstratos de Dados -- TAD}
    \Large

    Exemplos:

    \begin{itemize}
        \item TAD Ponto
        \item TAD Círculo
    \end{itemize}
\end{frame}

\begin{frame}[fragile]{Exercício: TAD Matriz}
    \begin{verbatim}
typedef struct matriz Matriz;

Matriz* mat_cria(int m, int n);
void mat_libera(Matriz* mat);
float mat_acessa(Matriz* mat, int i, int j);
void mat_atribui(Matriz* mat, int i, int j, float v);
int mat_linhas(Matriz* mat);
int mat_colunas(Matriz* mat);

struct matriz {
    int lin;
    int col;
    float* v;
    /* outra opção: matriz como vetor de ponteiros */
    /* float** v; */
}
    \end{verbatim}
\end{frame}

\begin{frame}[standout]{}
    \huge
    Dúvidas?
\end{frame}

\end{document}
